\chapter{Lebesgue Spaces}

Throughout this chapter, $\Omega\subseteq\mathbb R^d$ is measurable and all
$L^p$ spaces are taken with respect to Lebesgue measure unless stated
otherwise.  For $1<p<\infty$, the conjugate exponent is
$p'=p/(p-1)$.

\section{Uniform convexity, reflexivity, and duality}

\begin{notetheorem}
For $1<p<\infty$, the space $L^p(\Omega)$ is uniformly convex and therefore
reflexive.
\end{notetheorem}

\begin{proof}
Let $f,g\in L^p(\Omega)$ with $\|f\|_p,\|g\|_p\le1$ and
$\|f-g\|_p\ge\varepsilon$.  When $2\le p<\infty$, Clarkson's inequality is
\[
 \left\|\frac{f+g}{2}\right\|_p^p
 +\left\|\frac{f-g}{2}\right\|_p^p
 \le \frac{\|f\|_p^p+\|g\|_p^p}{2}\le1.
\]
Consequently,
\[
  \left\|\frac{f+g}{2}\right\|_p^p
  \le 1-\left(\frac{\varepsilon}{2}\right)^p,
\]
and hence
\[
  \left\|\frac{f+g}{2}\right\|_p
  \le 1-\delta_p(\varepsilon),
  \qquad
  \delta_p(\varepsilon)
  :=1-\left(1-(\varepsilon/2)^p\right)^{1/p}>0.
\]

For $1<p\le2$, put $q=p'=p/(p-1)$.  The second Clarkson inequality reads
\[
  \|f+g\|_p^q+\|f-g\|_p^q
  \le 2\bigl(\|f\|_p^p+\|g\|_p^p\bigr)^{q/p}.
\]
After division by $2^q$, the assumptions $\|f\|_p,\|g\|_p\le1$ give
\[
  \left\|\frac{f+g}{2}\right\|_p^q
  +\left\|\frac{f-g}{2}\right\|_p^q\le1.
\]
Thus the same argument, now with exponent $q$, gives
\[
  \left\|\frac{f+g}{2}\right\|_p
  \le \left(1-(\varepsilon/2)^q\right)^{1/q}<1.
\]
This is uniform convexity.  The Milman--Pettis theorem therefore implies that
$L^p(\Omega)$ is reflexive.
\end{proof}

\begin{notetheorem}[Riesz representation for $L^p$]
Let $1<p<\infty$.  For each $g\in L^{p'}(\Omega)$, the formula
\[
  T_g(f)=\int_\Omega f g\,\dd x
\]
defines an element of $(L^p(\Omega))^*$, and
\[
  \|T_g\|=\|g\|_{p'}.
\]
Every element of $(L^p(\Omega))^*$ is of this form.  Hence
\[
  (L^p(\Omega))^*\cong L^{p'}(\Omega)
\]
isometrically.
\end{notetheorem}

\begin{proof}
For $g\in L^{p'}(\Omega)$, Hölder's inequality gives
\[
  |T_g(f)|\le\|f\|_p\|g\|_{p'},
  \qquad f\in L^p(\Omega),
\]
so $\|T_g\|\le\|g\|_{p'}$.  If $g\ne0$, define
\[
  f_g(x)=
  \frac{|g(x)|^{p'-2}\overline{g(x)}}{\|g\|_{p'}^{p'-1}},
\]
with the value $0$ where $g=0$.  Since $(p'-1)p=p'$, one has
$\|f_g\|_p=1$, and
\[
  T_g(f_g)
  =\frac{1}{\|g\|_{p'}^{p'-1}}
    \int_\Omega |g|^{p'}\,\dd x
  =\|g\|_{p'}.
\]
Thus $\|T_g\|=\|g\|_{p'}$.

Let
\[
  J\colon L^{p'}(\Omega)\longrightarrow (L^p(\Omega))^*,
  \qquad Jg=T_g.
\]
The map $J$ is an isometry, so $J(L^{p'})$ is norm closed.  Suppose that
$\Phi\in(L^p)^{**}$ annihilates this range.  Since $L^p$ is reflexive, there is
$f\in L^p$ such that $\Phi=J_{L^p}f$.  Therefore
\[
  0=\Phi(Jg)=Jg(f)=\int_\Omega fg\,\dd x
  \qquad(g\in L^{p'}).
\]
Taking
\[
  g=|f|^{p-2}\overline f
\]
(with value $0$ at the zeros of $f$) is legitimate because
$|g|^{p'}=|f|^p$.  It gives
\[
  0=\int_\Omega |f|^p\,\dd x,
\]
so $f=0$.  Hence the annihilator of $J(L^{p'})$ in $(L^p)^{**}$ is trivial.
A proper closed linear subspace of a normed space has a nonzero annihilator by
Hahn--Banach; therefore $J(L^{p'})=(L^p)^*$.
\end{proof}

\begin{notetheorem}[The endpoint duality theorem]
For a $\sigma$-finite measure space, and in particular for Lebesgue measure on
$\Omega\subseteq\mathbb R^d$,
\[
  (L^1(\Omega))^*\cong L^\infty(\Omega)
\]
isometrically through the same pairing
$T_g(f)=\int_\Omega fg$.
\end{notetheorem}

\begin{proof}
For $g\in L^\infty(\Omega)$, Hölder's inequality gives
\[
  \left|\int_\Omega fg\,\dd x\right|
  \le\|f\|_1\|g\|_\infty,
\]
so $g\mapsto T_g$ is a linear contraction from $L^\infty$ into $(L^1)^*$.
We first prove surjectivity and then verify that the norm is preserved.

Let $T\in(L^1(\Omega))^*$.  Since the measure space is $\sigma$-finite, choose
measurable sets
\[
  E_1\subseteq E_2\subseteq\cdots,
  \qquad \mu(E_n)<\infty,
  \qquad \bigcup_{n=1}^\infty E_n=\Omega.
\]
For a measurable set $A\subseteq E_n$, define
\[
  \nu_n(A)=T(\one_A).
\]
If $(A_k)$ are pairwise disjoint subsets of $E_n$, then
\[
  \one_{\cup_{k=1}^m A_k}\longrightarrow
  \one_{\cup_{k=1}^\infty A_k}
  \quad\text{in }L^1
\]
because the omitted tail has measure tending to zero.  Continuity of $T$
therefore shows that $\nu_n$ is countably additive.  Moreover,
\[
  |\nu_n(A)|\le\|T\|\,\|\one_A\|_1
  =\|T\|\,\mu(A),
\]
so $\nu_n\ll\mu|_{E_n}$.  The Radon--Nikodym theorem gives a measurable
function $g_n$ on $E_n$ such that
\[
  \nu_n(A)=\int_A g_n\,\dd\mu
  \qquad(A\subseteq E_n\text{ measurable}).
\]
The displayed domination implies $|g_n|\le\|T\|$ almost everywhere.  Indeed,
if $|g_n|>\|T\|+\varepsilon$ on a set of positive measure, restricting to a
subset on which the complex phase of $g_n$ lies in a sufficiently short arc
(and using the sign in the real case) would contradict
$|\nu_n(A)|\le\|T\|\mu(A)$.

For $m>n$, the measures represented by $g_m$ and $g_n$ agree on every subset
of $E_n$.  Uniqueness in the Radon--Nikodym theorem gives
$g_m=g_n$ almost everywhere on $E_n$.  The functions therefore patch to a
measurable $g$ on $\Omega$ satisfying
\[
  \|g\|_\infty\le\|T\|.
\]
If $s$ is a simple function supported in some $E_n$, linearity gives
\[
  T(s)=\int_\Omega sg\,\dd\mu.
\]
Such simple functions are dense in $L^1(\Omega)$.  If $s_k\to f$ in $L^1$,
then
\[
  T(s_k)\to T(f)
  \quad\text{and}\quad
  \int s_kg\to\int fg,
\]
so $T(f)=\int fg$ for every $f\in L^1$.

It remains to prove equality of norms.  We already know
$\|T_g\|\le\|g\|_\infty$.  Put $M=\|g\|_\infty$ and let $\varepsilon>0$.
If $M>0$, the set $\{|g|>M-\varepsilon\}$ has positive measure.  Its
intersection with some $E_n$ contains a measurable subset $A$ with
$0<\mu(A)<\infty$.  Define
\[
  f=\frac{\one_A\,\overline{\operatorname{sgn}g}}{\mu(A)},
  \qquad
  \operatorname{sgn}g=
  \begin{cases}g/|g|,&g\ne0,\\0,&g=0.\end{cases}
\]
Then $\|f\|_1=1$ and
\[
  |T_g(f)|
  =\frac1{\mu(A)}\int_A|g|\,\dd\mu
  \ge M-\varepsilon.
\]
Letting $\varepsilon\downarrow0$ yields $\|T_g\|\ge M$.  Hence the
identification is isometric.
\end{proof}

\begin{noteremark}
Except in the finite-dimensional case (equivalently, when the measure algebra
has only finitely many atoms), $L^1(\Omega)$ is not reflexive.  Indeed, if it
were reflexive then its dual $L^\infty(\Omega)$ would also be reflexive, which
fails on every genuinely infinite measure algebra.  Likewise
$L^\infty(\Omega)$ is nonreflexive except in the same finite-dimensional
situation.
\end{noteremark}

\section{Separability and dense classes}

\begin{notetheorem}
For $1\le p<\infty$, the space $L^p(\mathbb R^d)$ is separable.  The space
$L^\infty(\mathbb R^d)$ is not separable.
\end{notetheorem}

\begin{proof}
Finite linear combinations, with rational coefficients, of indicators of
rectangles with rational endpoints form a countable set.  Simple-function
approximation and regularity of Lebesgue measure show that this set is dense in
$L^p$ for $p<\infty$.  For $L^\infty$, choose infinitely many pairwise
measurably distinguishable sets.  More explicitly, for $t\in\mathbb R$ let
\[
  A_t=(-\infty,t)\times\mathbb R^{d-1}.
\]
If $s\ne t$, then $A_s\triangle A_t$ has positive measure, and therefore
\[
  \|\one_{A_s}-\one_{A_t}\|_\infty=1.
\]
Thus $L^\infty(\mathbb R^d)$ contains an uncountable $1$-separated set and
cannot be separable.
\end{proof}

\begin{notetheorem}
For $1\le p<\infty$, $C_c(\mathbb R^d)$ and $C_c^\infty(\mathbb R^d)$ are
dense in $L^p(\mathbb R^d)$.
\end{notetheorem}

\begin{proof}
Fix $f\in L^p(\mathbb R^d)$ and $\varepsilon>0$.  Truncating first in space and
then in size gives
\[
  f_{R,M}=\one_{B(0,R)}\max\{-M,\min\{f,M\}\}
\]
in the real case, with the analogous radial truncation in the complex case,
such that $\|f-f_{R,M}\|_p<\varepsilon/3$ for large $R$ and $M$.  The bounded,
finite-support function $f_{R,M}$ can be approximated in $L^p$ by a simple
function
\[
  s=\sum_{j=1}^m a_j\one_{E_j},
\]
where every $E_j$ is measurable, bounded, and has finite measure.

It remains to approximate one indicator.  By inner and outer regularity of
Lebesgue measure, for any $\eta>0$ there are a compact set $K\subseteq E$ and
an open bounded set $O\supseteq E$ such that
\[
  |O\setminus K|<\eta.
\]
Choose $\phi\in C_c(\mathbb R^d)$ with
\[
  0\le\phi\le1,
  \qquad \phi=1\text{ on }K,
  \qquad \supp\phi\subseteq O.
\]
For instance, one may use the distance functions to $K$ and to
$\mathbb R^d\setminus O$.  Since $\one_E$ and $\phi$ can differ only on
$O\setminus K$,
\[
  \|\one_E-\phi\|_p^p\le |O\setminus K|<\eta.
\]
Choosing the errors for the finitely many $E_j$ sufficiently small produces
$v\in C_c(\mathbb R^d)$ with $\|s-v\|_p<\varepsilon/3$.  This proves density
of $C_c$.

Finally let $\rho_\delta$ be a mollifier.  If $v\in C_c$, then
$\rho_\delta*v\in C_c^\infty$, its support stays in a fixed compact set for
small $\delta$, and uniform continuity gives
\[
  \|\rho_\delta*v-v\|_\infty\longrightarrow0.
\]
The fixed finite support then gives convergence in $L^p$.  Combining the three
approximations proves density of $C_c^\infty$ as well.
\end{proof}

\section{Convolution and mollifiers}

\begin{notetheorem}[Young's convolution inequality]
Let $1\le p\le\infty$, $f\in L^1(\mathbb R^d)$, and
$g\in L^p(\mathbb R^d)$.  Then
\[
  (f*g)(x)=\int_{\mathbb R^d}f(x-y)g(y)\,\dd y
\]
is defined for almost every $x$, belongs to $L^p$, and satisfies
\[
  \|f*g\|_p\le\|f\|_1\|g\|_p.
\]
\end{notetheorem}

\begin{proof}
For $p=1$, Tonelli's theorem and translation invariance give
\[
\begin{aligned}
  \|f*g\|_1
  &\le\int_{\mathbb R^d}\int_{\mathbb R^d}
       |f(x-y)|\,|g(y)|\,\dd y\,\dd x\\
  &=\|f\|_1\|g\|_1.
\end{aligned}
\]
For $p=\infty$, for almost every $x$,
\[
  |(f*g)(x)|
  \le\int |f(x-y)|\,|g(y)|\,\dd y
  \le\|g\|_\infty\|f\|_1.
\]
For $1<p<\infty$, Hölder's inequality with respect to the measure
$|f(x-y)|\dd y$ gives
\[
 |(f*g)(x)|^p
 \le \|f\|_1^{p-1}
     \int |f(x-y)|\,|g(y)|^p\,\dd y.
\]
Integrating in $x$ and applying Tonelli yields the result.
\end{proof}

\begin{noteproposition}[Differentiation of convolution]
If $f\in C_c^k(\mathbb R^d)$ and $g\in L^1_{\mathrm{loc}}(\mathbb R^d)$,
then $f*g\in C^k$ wherever the convolution is defined, and
\[
  D^\alpha(f*g)=(D^\alpha f)*g
  \qquad(|\alpha|\le k).
\]
\end{noteproposition}

\begin{proof}
Fix a compact set $K\subset\mathbb R^d$.  If $x\in K$ and
$f(x-y)\ne0$, then
\[
  y\in x-\supp f\subseteq K-\supp f,
\]
which is compact.  Thus all relevant values of $g$ lie in one fixed compact
set on which $g$ is integrable.  For a coordinate vector $e_j$,
\[
  \frac{(f*g)(x+he_j)-(f*g)(x)}{h}
  =\int_{\mathbb R^d}
     \frac{f(x+he_j-y)-f(x-y)}{h}\,g(y)\,\dd y.
\]
For small $h$, the difference quotient is supported in one fixed compact set
and is bounded there by a constant multiple of $|g(y)|$, by the mean-value
theorem.  Dominated convergence therefore gives
\[
  \partial_j(f*g)(x)=((\partial_jf)*g)(x).
\]
The same domination, now applied to differences of the last integral, gives
continuity in $x$.  Iterating the argument proves the formula for every
$|\alpha|\le k$.
\end{proof}

\begin{notedefinition}
A sequence $(\rho_n)$ is an \emph{approximate identity} or a sequence of
\emph{mollifiers} if
\[
  \rho_n\in C_c^\infty(\mathbb R^d),\qquad
  \rho_n\ge0,
  \qquad
  \int\rho_n=1,
  \qquad
  \supp\rho_n\subseteq B(0,1/n).
\]
A standard choice is $\rho_n(x)=n^d\rho(nx)$ for a fixed nonnegative
$\rho\in C_c^\infty(B(0,1))$ with integral one.
\end{notedefinition}

\begin{noteproposition}
Let $(\rho_n)$ be mollifiers.
\begin{enumerate}
  \item If $f\in C_c(\mathbb R^d)$, then $\rho_n*f\to f$ uniformly.
  \item If $f\in L^p(\mathbb R^d)$ with $1\le p<\infty$, then
        $\rho_n*f\to f$ in $L^p$.
\end{enumerate}
\end{noteproposition}

\begin{proof}
For continuous compactly supported $f$,
\[
 |(\rho_n*f)(x)-f(x)|
 \le\int\rho_n(y)|f(x-y)-f(x)|\,\dd y,
\]
and uniform continuity makes the right-hand side tend uniformly to zero.  The
$L^p$ statement follows by approximating $f$ by a function in $C_c$, using
Young's inequality for the two error terms.
\end{proof}

\begin{notecorollary}[Vanishing weak gradient]
Let $\Omega\subseteq\mathbb R^d$ be open and connected, and let
$u\in L^1_{\mathrm{loc}}(\Omega)$.  If
\[
  \int_\Omega u\,\partial_j\varphi\,\dd x=0
  \qquad
  (\varphi\in C_c^\infty(\Omega),\ 1\le j\le d),
\]
then $u$ is almost everywhere equal to a constant on $\Omega$.
\end{notecorollary}

\begin{proof}
Fix a ball $B(x_0,2r)\Subset\Omega$ and extend $u$ arbitrarily outside
$B(x_0,2r)$ for the purpose of convolution.  For
$0<\varepsilon<r$, put $u_\varepsilon=\rho_\varepsilon*u$ on
$B(x_0,r)$.  For $1\le j\le d$ and $x\in B(x_0,r)$,
\[
  \partial_j u_\varepsilon(x)
  =\int u(y)\,\partial_j\rho_\varepsilon(x-y)\,\dd y=0,
\]
because $y\mapsto\rho_\varepsilon(x-y)$ is a test function supported in
$B(x_0,2r)$.  Hence $u_\varepsilon$ is constant on the connected ball
$B(x_0,r)$; write this constant as $c_\varepsilon$.

Choose a smaller ball $B_0\Subset B(x_0,r)$ of positive measure.  Since
$u_\varepsilon\to u$ in $L^1(B_0)$, the constants $(c_\varepsilon)$ form a
Cauchy net:
\[
  |B_0|\,|c_\varepsilon-c_\delta|
  =\|c_\varepsilon-c_\delta\|_{L^1(B_0)}
  \le\|u_\varepsilon-u\|_{L^1(B_0)}
     +\|u-u_\delta\|_{L^1(B_0)}.
\]
Thus $c_\varepsilon\to c$, and $u=c$ almost everywhere on $B_0$, hence on
$B(x_0,r)$ by the same $L^1$ convergence.  We have proved local constancy on
every compactly contained ball.  Along a finite chain of overlapping balls,
the constants agree on intersections of positive measure.  Since an open
connected subset of $\mathbb R^d$ is path connected by polygonal chains of
balls, one constant propagates throughout $\Omega$.
\end{proof}
